\documentclass[11pt]{article}
\usepackage[margin=1in]{geometry}
\usepackage{amsmath,amssymb,amsthm}
\usepackage[utf8]{inputenc}
\usepackage{listings}
\usepackage{hyperref}
\usepackage{xcolor}
\usepackage{textcomp}

\definecolor{leanblue}{rgb}{0.2,0.4,0.7}

\lstdefinelanguage{Lean}{
  keywords={theorem,lemma,def,axiom,structure,class,instance,where,by,sorry,simp,exact,intro,apply,have,let,fun,match,if,then,else,import,open,namespace,end,section,variable,inductive,noncomputable,abbrev,deriving},
  keywordstyle=\color{leanblue}\bfseries,
  comment=[l]{--},
  morecomment=[s]{/-}{-/},
  string=[b]",
  basicstyle=\ttfamily\small,
  breaklines=true,
  extendedchars=true,
  inputencoding=utf8,
  literate=
    {≥}{{$\geq$}}1
    {≤}{{$\leq$}}1
    {→}{{$\to$}}1
    {←}{{$\leftarrow$}}1
    {∧}{{$\wedge$}}1
    {∨}{{$\vee$}}1
    {¬}{{$\neg$}}1
    {∀}{{$\forall$}}1
    {∃}{{$\exists$}}1
    {⊆}{{$\subseteq$}}1
    {⊂}{{$\subset$}}1
    {∈}{{$\in$}}1
    {∉}{{$\notin$}}1
    {×}{{$\times$}}1
    {⊗}{{$\otimes$}}1
    {▸}{{$\blacktriangleright$}}1
    {⟨}{{$\langle$}}1
    {⟩}{{$\rangle$}}1
    {λ}{{$\lambda$}}1
    {α}{{$\alpha$}}1
    {β}{{$\beta$}}1
    {σ}{{$\sigma$}}1
    {θ}{{$\theta$}}1
    {ε}{{$\varepsilon$}}1
    {μ}{{$\mu$}}1
    {π}{{$\pi$}}1
    {Σ}{{$\Sigma$}}1
    {Π}{{$\Pi$}}1
    {▶}{{$\triangleright$}}1
    {⊔}{{$\sqcup$}}1
    {⊓}{{$\sqcap$}}1
    {⊥}{{$\bot$}}1
    {⊤}{{$\top$}}1
    {∅}{{$\emptyset$}}1
    {∪}{{$\cup$}}1
    {∩}{{$\cap$}}1
    {≠}{{$\neq$}}1
    {ℕ}{{$\mathbb{N}$}}1
    {₀}{{$_0$}}1
    {₁}{{$_1$}}1
    {₂}{{$_2$}}1
    {|>}{{$|{\!>}$}}2
    {·}{{$\cdot$}}1
    {—}{{---}}1
    {∘}{{$\circ$}}1
    {√}{{$\sqrt{}$}}1
    {∝}{{$\propto$}}1
    {≈}{{$\approx$}}1
    {═}{{=}}1
    {⟹}{{$\Longrightarrow$}}1,
}

\newtheorem{theorem}{Theorem}
\newtheorem{lemma}[theorem]{Lemma}
\newtheorem{definition}[theorem]{Definition}
\newtheorem{axiom_stmt}[theorem]{Axiom}

\title{Formal Verification: Consensus Safety}
\author{Zach Kelling, Lux Industries}
\date{December 2025}

\begin{document}
\maketitle

\begin{abstract}
We prove that no two honest nodes finalize conflicting values at the same height in the Lux consensus protocol. The proof models the consensus as a state machine with preferences, confidence counters, and finalization thresholds. The main safety theorem is proved from two axioms: unique threshold (at most one value achieves $\alpha$-quorum per round) and finalization window overlap (two finalized values share a common round).
\end{abstract}

\section{Introduction}
This is the core safety theorem for Lux consensus. The proof follows the standard metastable BFT argument: (1) honest majority ensures at most one value wins per round, (2) consecutive wins accumulate confidence, (3) once finalized, preference is locked, (4) therefore two honest finalizations must agree. The model maps directly to \texttt{protocol/wave/wave.go} (sampling, threshold, confidence) and \texttt{core/consensus.go} (Block interface).

\section{Definitions}
\begin{definition}[ValueId]
A value identifier modeled as $\text{Fin}(2^{256})$, corresponding to \texttt{ids.ID = [32]byte} in Go.
\end{definition}

\begin{definition}[Params]
Consensus parameters: $n$ nodes, $f$ max Byzantine, $k$ sample size, $\alpha$ vote threshold, $\beta$ confidence threshold, with invariants $f < n/3$, $k \leq n$, $\alpha \leq k$, $0 < \beta$.
\end{definition}

\begin{definition}[NodeState]
Per-node state: preference, confidence counter, finalized value, honesty flag.
\end{definition}

\begin{definition}[ConsensusState]
Global state: round number and per-node states.
\end{definition}


\section{Theorems and Axioms}
\begin{axiom_stmt}[\texttt{finalized\_preference\_stable}]
Once a node finalizes value $v$, it remains finalized in all future rounds. Models the early-return guard in \texttt{wave.go}.
\end{axiom_stmt}

\begin{axiom_stmt}[\texttt{unique\_threshold}]
If honest count $> 2f$ and $\alpha > 2k/3$, then at most one value can achieve $\alpha$-threshold preference among honest nodes.
\end{axiom_stmt}

\begin{axiom_stmt}[\texttt{finalization\_windows\_overlap}]
If two nodes finalize (possibly different) values, their $\beta$-length confirmation windows share at least one common round.
\end{axiom_stmt}

\begin{theorem}[\texttt{safety}]
If two honest nodes $i, j$ finalize values $v_1, v_2$ respectively, and the BFT assumption holds ($\text{honestCount} > 2f$, $\alpha > 2k/3$), then $v_1 = v_2$.
\end{theorem}
\begin{proof}
From \texttt{finalization\_windows\_overlap}, obtain a state $s_{\text{overlap}}$ where both values achieved $\alpha$-threshold. Apply \texttt{unique\_threshold} to conclude $v_1 = v_2$.
\end{proof}


\section{Lean Source}
\lstinputlisting[language=Lean]{../../formal/lean/Consensus/Safety.lean}

\section{References}
\noindent Source code: \url{https://github.com/luxfi/proofs/blob/main/lean/Consensus/Safety.lean}\\
\noindent White papers: \url{https://papers.lux.network}

\noindent Related white papers:
\begin{itemize}
  \item \texttt{lux-consensus.tex} --- Lux Consensus Architecture
\end{itemize}

\end{document}
